\bigskip\Needspace{9\baselineskip}
\pdfbookmark[1]{Ceres Ferdinandea. 1802 Oktober}{ceres180210}
\begin{center}
CERES FERDINANDEA.\par
\medskip\hrule\medskip
Monatliche Correspondenz. Band VI. S.\ 382—389. 1802 Oktober.\par
\medskip\hrule
\end{center}\medskip
—\quad —\quad —

Mehrere Astronomen haben die Störungen der Länge und des Abstands der \textit{Ceres} durch Jupiter bereits berechnet. Solche verwickelten und schwierigen Rechnungen kann man jedoch gar nicht oft genug wiederholen. Erst kürzlich schrieb uns Prof.\ \textsc{Wurm} darüber, als er uns einige Verbesserungen und erläuternde Ergänzungen zu den Formeln für die
\newcommand{\JupiterSymbol}{\text{\mdseries\AstroSymbols ♃}}
\SourceLeaf{230}%
Störung durch Mars schickte, die wir demnächst in unseren Heften mitteilen werden: Ich glaube überhaupt, wer auf diesem Weg der Störungsrechnungen noch nie auf Dornen gestoßen ist, hat ihn bestimmt nie betreten. Denn bei solchen Untersuchungen lässt sich wohl kaum vermeiden, dass man hin und wieder in die Dornen gerät und vielleicht sogar, bis man sich besser zurechtfindet, ein kleines \textit{quid pro quo} begeht.

Dr.\ \textsc{Gauß} hat diese Störungsrechnungen daher mit seinen Elementen VII der Ceresbahn wiederholt und hier und da kleine Unterschiede gefunden. Die analytischen Formeln, die er bei dieser Rechnung verwendet, hat er zunächst alle selbst entwickelt. Einige davon haben eine etwas andere Form als die von \textsc{La Place}.
\begin{center}
\textit{Störungen der Ceres Ferdinandea durch Jupiter, berechnet von Dr.\ }\textsc{Gauß}\textit{.}\par\medskip
{\small\setlength{\tabcolsep}{3pt}
\begin{tabular}{l r}
Jährliche Abnahme der Exzentrizität $\ .\ .\ .\ .\ .\ .\ .\ .\ .\ .\ .\ .$&$0.000005909$\\
Jährliche Bewegung des Aphels gegenüber den Fixsternen $\ .\ .\ .\ .\ .$&$+70''15$
\end{tabular}}
\par\Needspace{14\baselineskip}\bigskip
\textit{Periodische Gleichungen.}\par\medskip
1)\quad \textbf{Für die Länge in der Bahn.}
\end{center}
\noindent\begin{minipage}[t]{.46\linewidth}\vspace{0pt}\small\setlength{\arraycolsep}{2pt}
\[
\begin{array}{l}
-231''94\sin(\CeresSymbol-\JupiterSymbol)\\
+496.68\sin2(\CeresSymbol-\JupiterSymbol)\\
+44.15\sin3(\CeresSymbol-\JupiterSymbol)\\
+10.07\sin4(\CeresSymbol-\JupiterSymbol)\\
+3.05\sin5(\CeresSymbol-\JupiterSymbol)\\
+1.07\sin6(\CeresSymbol-\JupiterSymbol)\\
+0.41\sin7(\CeresSymbol-\JupiterSymbol)\\
-60.25\sin(\JupiterSymbol+17^0\ 41'\ 28'')\\
-621.29\sin(\CeresSymbol-2\JupiterSymbol-26^0\ 49'\ 22'')
\end{array}
\]
\end{minipage}\hfill\vrule width .35pt\hfill\begin{minipage}[t]{.515\linewidth}\vspace{0pt}\small\setlength{\arraycolsep}{2pt}
\[
\begin{array}{l}
-442''65\sin(2\CeresSymbol-3\JupiterSymbol-11^0\ 34'\ 39'')\\
+56.69\sin(3\CeresSymbol-4\JupiterSymbol-11\quad24\quad1'')\\
+10.58\sin(4\CeresSymbol-5\JupiterSymbol-11\quad11\quad21'')\\
+3.26\sin(5\CeresSymbol-6\JupiterSymbol-10\quad55\quad43'')\\
+23.62\sin(2\CeresSymbol-\JupiterSymbol+36\quad31\quad34'')\\
-53.93\sin(3\CeresSymbol-2\JupiterSymbol+33\quad24\quad19'')\\
-5.96\sin(4\CeresSymbol-3\JupiterSymbol+31\quad20\quad18'')\\
-1.70\sin(5\CeresSymbol-4\JupiterSymbol+29\quad33\quad57'')
\end{array}
\]
\end{minipage}\par
\Needspace{15\baselineskip}
\begin{center}2)\quad \textbf{Für den Radiusvektor.}\end{center}
\noindent\begin{minipage}[t]{.46\linewidth}\vspace{0pt}\small\setlength{\arraycolsep}{2pt}
\[
\begin{array}{l}
-0.0000947\\
+0.0010304\cos(\CeresSymbol-\JupiterSymbol)\\
-0.0038023\cos2(\CeresSymbol-\JupiterSymbol)\\
-0.0004206\cos3(\CeresSymbol-\JupiterSymbol)\\
-0.0001077\cos4(\CeresSymbol-\JupiterSymbol)\\
-0.0000350\cos5(\CeresSymbol-\JupiterSymbol)\\
-0.0000128\cos6(\CeresSymbol-\JupiterSymbol)\\
-0.0000050\cos7(\CeresSymbol-\JupiterSymbol)\\
-0.0000633\cos(\CeresSymbol+25^0\ 56'\ 46'')\\
+0.0002469\cos(\JupiterSymbol+23\quad42\quad53'')
\end{array}
\]
\end{minipage}\hfill\vrule width .35pt\hfill\begin{minipage}[t]{.515\linewidth}\vspace{0pt}\small\setlength{\arraycolsep}{2pt}
\[
\begin{array}{l}
+0.0008589\cos(\CeresSymbol-2\JupiterSymbol-24^0\ 37'\ 54'')\\
+0.0025827\cos(2\CeresSymbol-3\JupiterSymbol-11\quad19\quad13'')\\
-0.0004817\cos(3\CeresSymbol-4\JupiterSymbol-11\quad32\quad33'')\\
-0.0001075\cos(4\CeresSymbol-5\JupiterSymbol-11\quad31\quad15'')\\
-0.0000366\cos(5\CeresSymbol-6\JupiterSymbol-11\quad21\quad15'')\\
-0.0001367\cos(2\CeresSymbol-\JupiterSymbol+37\quad40\quad22'')\\
+0.0003049\cos(3\CeresSymbol-2\JupiterSymbol+32\quad54\quad59'')\\
+0.0000411\cos(4\CeresSymbol-3\JupiterSymbol+29\quad43\quad48'')\\
+0.0000139\cos(5\CeresSymbol-4\JupiterSymbol+27\quad40\quad24'')
\end{array}
\]
\end{minipage}\par
\begin{flushright}[\textsc{von Zach}.]\end{flushright}
